\begin{table}[htbp]
\centering
\small
\setlength{\tabcolsep}{3pt}
\caption{Mathematical techniques used in continuum wave analysis.}
\label{tab:T5-techniques}
\begin{tabularx}{\textwidth}{@{}p{2.4cm}YYYY@{}}
\toprule
\textbf{Technique} & \textbf{Mathematical basis} & \textbf{Typical use} & \textbf{Model-specific consideration} & \textbf{\shortstack{Reference/\\example}} \\
\midrule
Normal mode expansion & Separation of variables; eigenfunction expansion & Layered media and separable geometries & Boundary conditions determine admissible modes & Standard method \\
Green's functions & Fundamental solution and convolution & Point sources and infinite domains & Anisotropy may complicate integral evaluation & Standard method \\
Fourier transform & Frequency-domain representation & Time-harmonic problems and dispersion & Initial-value problems require time-domain treatment & Standard method \\
Laplace transform & Complex-frequency representation & Initial-value problems and transient waves & Inversion may require branch-cut treatment & Standard method \\
Christoffel equation & $\det[\Gamma_{ik}-\rho c^2\delta_{ik}]=0$ & Plane waves in anisotropic media & Polynomial order depends on the field equations & Standard method \\
Scattering matrix & Interface matching formulation & Coupled interface problems & Dimensions and normalization depend on the selected fields & Standard method \\
Finite element & Weak form and variational principles & Complex geometries and nonlinear problems & Mesh choice and computational cost & Standard method \\
\bottomrule
\end{tabularx}
\par\smallskip{\footnotesize These are general method roles; convergence and computational behavior depend on the selected problem, discretization, and implementation.}
\end{table}
